\clearpage
\section{TX-side Compression and Safety --- Upper Sideband}
\label{sec:compression-upper}

\textbf{What's plotted, and why it matters.} This section shows how the unit's transmitter responds as you
drive it harder. At 9 representative (FR1, FR3) points, the input power was swept upward from −60\,dBm until the gain had
dropped 5\,dB below its small-signal value (or the 0\,dBm limit of the test source was reached), with the same ``optimal''
filter-selection algorithm described in Section~\ref{sec:sideband-suppression}.\footnote{Every point on every curve is a real measurement,
never a simulation or an interpolation between sparse samples.} At low input power, the output just scales
linearly with the input --- this is the small-signal gain, shown as the dashed line and labeled directly on
each panel. Push the input harder, and the amplifiers can't keep up: the real curve bends away below that line.
This is called compression. P1dB marks where the gain has dropped 1\,dB below its small-signal value; P3dB
marks where it has dropped 3\,dB. Driving past P1dB starts to distort your signal, and driving well past P3dB
risks damaging the unit. Together, these two points tell you the safe input power range for this unit ---
knowing them is essential before connecting any real transmit signal.

\textbf{Measurement conditions.} \lmxDriveRule{} (set from the LO frequency), with the LO leakage suppressed by the unit's stored LO-null calibration (Section~\ref{sec:lo-suppression}), exactly as in normal use. The small-signal gain is taken from the readings between −50 and −46\,dBm at the source, because the test source is less accurate in its lowest few dB. The unit was measured hot, as the worst case: no fan, LMX2820 die at \condCompMinC--\condCompMaxC\textdegree C. A cooler unit will typically do somewhat better.

\textbf{Gain expansion before compression.} For some configurations, you'll notice the measured gain briefly
\emph{rises} above its small-signal value before it eventually compresses --- the curve bulges above the dashed
line rather than only bending below it.\footnote{This is normal, expected behavior for this unit, not a
measurement error. The board's final-stage power amplifier (Analog Devices ADL9006) is a cascode distributed
pHEMT amplifier, a class well known for exactly this kind of non-monotonic gain-vs-drive behavior. You can see
the same effect directly in the ADL9006's own published datasheet (Figures 21, 22, 24, and 25: P\textsubscript{out},
Gain, PAE, and I\textsubscript{DD} vs.\ input power, at 2, 10, 20, and 26\,GHz), where the gain trace rises before it
falls. Most introductory textbooks only describe gain compression and skip gain expansion entirely --- so if
this is the first time you've seen it, that is expected too.} You can see the same effect directly in the
ADL9006's own published datasheet.

\figpair{images/sec4_fig1.pdf}{FR1 = 1.0\,GHz, FR3 = 9.0\,GHz. Filter: 0x3E.}{fig:4-1}%
        {images/sec4_fig2.pdf}{FR1 = 2.0\,GHz, FR3 = 9.0\,GHz. Filter: 0x3D.}{fig:4-2}
\figpair{images/sec4_fig3.pdf}{FR1 = 3.5\,GHz, FR3 = 11.5\,GHz. Filter: 0x3E.}{fig:4-3}%
        {images/sec4_fig4.pdf}{FR1 = 1.0\,GHz, FR3 = 13.0\,GHz. Filter: 0x1C.}{fig:4-4}
\figpair{images/sec4_fig5.pdf}{FR1 = 2.0\,GHz, FR3 = 15.0\,GHz. Filter: 0x1E.}{fig:4-5}%
        {images/sec4_fig6.pdf}{FR1 = 3.0\,GHz, FR3 = 17.0\,GHz. Filter: 0x22.}{fig:4-6}

\figpair{images/sec4_fig7.pdf}{FR1 = 1.0\,GHz, FR3 = 19.0\,GHz. Filter: 0x2B.}{fig:4-7}%
        {images/sec4_fig8.pdf}{FR1 = 2.0\,GHz, FR3 = 22.0\,GHz. Filter: 0x2E.}{fig:4-8}
\figonesmall{images/sec4_fig9.pdf}{FR1 = 3.5\,GHz, FR3 = 23.5\,GHz. Filter: 0x2F. Tip: this configuration compresses early (output P1dB +8.1\,dBm, against +14 to +18\,dBm at most other points). Near the top of the IF band (FR1~=~3.5\,GHz) with a high LO, TX compression starts earlier; if you need more output power, use a lower FR1 for the same FR3. Running the unit cooler helps only a little: with the LMX2820 die at about 60\textdegree C instead of 83\textdegree C, the output P1dB here rose by about 3\,dB (from +8.1 to +11.3\,dBm).}{fig:4-9}
